% pandoc 模板:版式与 assignment01/handout/assignment01.tex 相同,
% 正文由 handout/src/*.md 生成,不要手工编辑生成出来的 .tex。
\documentclass[11pt,a4paper,fontset=fandol]{ctexart}

\usepackage{geometry}
\geometry{margin=2.3cm}

\usepackage{amsmath,amssymb}
\usepackage{booktabs,array,multirow,longtable}
\usepackage{calc,etoolbox}
\usepackage{enumitem}
\setlist{nosep,leftmargin=2em}
\usepackage{xcolor}
\definecolor{accent}{RGB}{28,60,120}
\definecolor{codebg}{RGB}{248,248,246}
\usepackage[most]{tcolorbox}
\usepackage{listings}
\definecolor{strgreen}{RGB}{21,101,42}
\lstset{
  basicstyle=\ttfamily\small,
  backgroundcolor=\color{codebg},
  frame=single, rulecolor=\color{gray!40}, framerule=0.4pt,
  breaklines=true, columns=fullflexible, keepspaces=true,
  showstringspaces=false,
  language=bash,
  morekeywords={make,nvcc,uv,pytest,python,cd,ncu,nvdisasm},
  commentstyle=\color{gray!130}\itshape,
  keywordstyle=\color{accent}\bfseries,
  stringstyle=\color{strgreen},
  xleftmargin=4pt, xrightmargin=4pt,
}
\usepackage{hyperref}
\hypersetup{colorlinks, linkcolor=accent, urlcolor=accent}
\emergencystretch=3em
\pagestyle{plain}

% pandoc 产物需要的宏
\providecommand{\tightlist}{\setlength{\itemsep}{0pt}\setlength{\parskip}{0pt}}
\providecommand{\passthrough}[1]{#1}
\setlength{\LTpre}{0.6em}
\setlength{\LTpost}{0.6em}

\usepackage{color}
\usepackage{fancyvrb}
\newcommand{\VerbBar}{|}
\newcommand{\VERB}{\Verb[commandchars=\\\{\}]}
\DefineVerbatimEnvironment{Highlighting}{Verbatim}{commandchars=\\\{\}}
% Add ',fontsize=\small' for more characters per line
\newenvironment{Shaded}{}{}
\newcommand{\AlertTok}[1]{\textcolor[rgb]{1.00,0.00,0.00}{\textbf{#1}}}
\newcommand{\AnnotationTok}[1]{\textcolor[rgb]{0.38,0.63,0.69}{\textbf{\textit{#1}}}}
\newcommand{\AttributeTok}[1]{\textcolor[rgb]{0.49,0.56,0.16}{#1}}
\newcommand{\BaseNTok}[1]{\textcolor[rgb]{0.25,0.63,0.44}{#1}}
\newcommand{\BuiltInTok}[1]{\textcolor[rgb]{0.00,0.50,0.00}{#1}}
\newcommand{\CharTok}[1]{\textcolor[rgb]{0.25,0.44,0.63}{#1}}
\newcommand{\CommentTok}[1]{\textcolor[rgb]{0.38,0.63,0.69}{\textit{#1}}}
\newcommand{\CommentVarTok}[1]{\textcolor[rgb]{0.38,0.63,0.69}{\textbf{\textit{#1}}}}
\newcommand{\ConstantTok}[1]{\textcolor[rgb]{0.53,0.00,0.00}{#1}}
\newcommand{\ControlFlowTok}[1]{\textcolor[rgb]{0.00,0.44,0.13}{\textbf{#1}}}
\newcommand{\DataTypeTok}[1]{\textcolor[rgb]{0.56,0.13,0.00}{#1}}
\newcommand{\DecValTok}[1]{\textcolor[rgb]{0.25,0.63,0.44}{#1}}
\newcommand{\DocumentationTok}[1]{\textcolor[rgb]{0.73,0.13,0.13}{\textit{#1}}}
\newcommand{\ErrorTok}[1]{\textcolor[rgb]{1.00,0.00,0.00}{\textbf{#1}}}
\newcommand{\ExtensionTok}[1]{#1}
\newcommand{\FloatTok}[1]{\textcolor[rgb]{0.25,0.63,0.44}{#1}}
\newcommand{\FunctionTok}[1]{\textcolor[rgb]{0.02,0.16,0.49}{#1}}
\newcommand{\ImportTok}[1]{\textcolor[rgb]{0.00,0.50,0.00}{\textbf{#1}}}
\newcommand{\InformationTok}[1]{\textcolor[rgb]{0.38,0.63,0.69}{\textbf{\textit{#1}}}}
\newcommand{\KeywordTok}[1]{\textcolor[rgb]{0.00,0.44,0.13}{\textbf{#1}}}
\newcommand{\NormalTok}[1]{#1}
\newcommand{\OperatorTok}[1]{\textcolor[rgb]{0.40,0.40,0.40}{#1}}
\newcommand{\OtherTok}[1]{\textcolor[rgb]{0.00,0.44,0.13}{#1}}
\newcommand{\PreprocessorTok}[1]{\textcolor[rgb]{0.74,0.48,0.00}{#1}}
\newcommand{\RegionMarkerTok}[1]{#1}
\newcommand{\SpecialCharTok}[1]{\textcolor[rgb]{0.25,0.44,0.63}{#1}}
\newcommand{\SpecialStringTok}[1]{\textcolor[rgb]{0.73,0.40,0.53}{#1}}
\newcommand{\StringTok}[1]{\textcolor[rgb]{0.25,0.44,0.63}{#1}}
\newcommand{\VariableTok}[1]{\textcolor[rgb]{0.10,0.09,0.49}{#1}}
\newcommand{\VerbatimStringTok}[1]{\textcolor[rgb]{0.25,0.44,0.63}{#1}}
\newcommand{\WarningTok}[1]{\textcolor[rgb]{0.38,0.63,0.69}{\textbf{\textit{#1}}}}

% 参考资料框
\newtcolorbox{reading}{enhanced, breakable,
  colback=gray!5, colframe=gray!50, boxrule=0.5pt, arc=1.5pt,
  left=7pt, right=7pt, top=5pt, bottom=5pt,
  title={参考资料}, fonttitle=\bfseries\sffamily\small,
  coltitle=black, attach title to upper={\quad}}

% 压轴题框
\newtcolorbox{capstone}[1]{enhanced, breakable,
  colback=accent!4, colframe=accent!70, boxrule=0.7pt, arc=1.5pt,
  left=7pt, right=7pt, top=5pt, bottom=5pt,
  title={#1}, fonttitle=\bfseries\sffamily}

% 回望框
\newtcolorbox{lookback}{enhanced, breakable,
  colback=gray!5, colframe=gray!50, boxrule=0.5pt, arc=1.5pt,
  left=7pt, right=7pt, top=5pt, bottom=5pt,
  title={Editor's Note}, fonttitle=\bfseries\sffamily\small,
  coltitle=black, attach title to upper={\quad},
  before upper app={\setlength{\parskip}{0.5em}}}

% 学员作答
\newtcolorbox{answer}{enhanced, breakable,
  colback=strgreen!3, colframe=strgreen!55, boxrule=0.5pt, arc=1.5pt,
  left=7pt, right=7pt, top=5pt, bottom=5pt,
  title={作答}, fonttitle=\bfseries\sffamily\small,
  coltitle=black, attach title to upper={\quad},
  before upper app={\setlength{\parskip}{0.4em}}}

% 题目标题:\prob{编号}{TYPE};选做 \prob[Optional]{编号}{TYPE};
% 末尾可选参数标注涉及的代码文件
\NewDocumentCommand{\prob}{O{}mmo}{\par\medskip\noindent
  {\bfseries prob #2}\;{\small\sffamily(#3\IfBlankF{#1}{\,·\,#1})}%
  \IfValueT{#4}{\hfill{\footnotesize\ttfamily\color{black!60}#4}}\par\nobreak\smallskip\noindent\ignorespaces}

\begin{document}

\begin{center}
  {\LARGE\bfseries 作业 2:Tensor Core \& Pipeline}\\[6pt]
  {\sffamily Weiming HPC Training Camp \(\times\) LCPU AI Infra Seminars
· Session 3 / 4}
\end{center}

\vspace{0.5em}
\hrule
\vspace{1em}

\setcounter{section}{-1}

\section*{Preface}\label{preface}
\addcontentsline{toc}{section}{Preface}

这次作业主要是 Session 3 的配套练习，内容围绕 Tensor Core 展开，其中
Module 4 还会用到 Session 4 介绍的 tiling、TMA 和
pipeline。完成作业时主要参考课程课件和 PTX ISA 文档。

题型沿用系列惯例，并新增 DERIVE：先手工推导，再写程序验证推导。共七种：
CONCEPT 概念题、DERIVE 推导题、MODIFY 改造题、DEBUG 修 bug 题、
EXPERIMENT 实验题、HANDS-ON 动手题、FROM-SCRATCH 从零实现题。标 Optional
的为选做。涉及代码的题在标题行右侧标出文件路径(相对
\texttt{assignment02/})，题面细节以文件头注释为准；FROM-SCRATCH 题。

硬件与构建:主线环境是 B300，\texttt{cuda/Makefile} 默认
\texttt{ARCH=100f}；M0、M1 也可以在 5090
上完成(\texttt{ARCH=120a\ make\ ...})； M2 的判测纯
host，无卡可判；M3、M4、M5 需要 B300。本作业统一用显式 \texttt{-gencode}
而不用 \texttt{-arch=sm\_XXXa} 简写，原因见
\texttt{assignment02/README.md}，Makefile 已经配好。

实验纪律:所有性能数字都在自己占满的卡上测；测前测后
\texttt{nvidia-smi\ -\/-query-compute-apps=pid，name\ -\/-format=csv}
查有没有别的 进程占卡；可能 hang 的程序套
\texttt{timeout\ -k\ 5\ \textless{}秒数\textgreater{}}(判测脚本已带)。

AI政策:必做题沿用仓库根目录 \texttt{CLAUDE.md}------AI 可以帮你理解、
review，不能替你实现；团队题不设限制。详见
\texttt{assignment02/README.md}。

\subsection*{Content}\label{content}
\addcontentsline{toc}{subsection}{Content}

\begin{longtable}[]{@{}
  >{\raggedright\arraybackslash}p{(\linewidth - 6\tabcolsep) * \real{0.2500}}
  >{\raggedright\arraybackslash}p{(\linewidth - 6\tabcolsep) * \real{0.2500}}
  >{\raggedright\arraybackslash}p{(\linewidth - 6\tabcolsep) * \real{0.2500}}
  >{\raggedright\arraybackslash}p{(\linewidth - 6\tabcolsep) * \real{0.2500}}@{}}
\toprule\noalign{}
\begin{minipage}[b]{\linewidth}\raggedright
Module
\end{minipage} & \begin{minipage}[b]{\linewidth}\raggedright
主题
\end{minipage} & \begin{minipage}[b]{\linewidth}\raggedright
对应课件
\end{minipage} & \begin{minipage}[b]{\linewidth}\raggedright
代码位置
\end{minipage} \\
\midrule\noalign{}
\endhead
\bottomrule\noalign{}
\endlastfoot
0 & 环境与峰值 & P1(S008--S021) & \texttt{cuda/m0\_env/} \\
1 & sm80:fragment 与 mma.sync & P2.2(S025--S041) &
\texttt{cuda/m1\_sm80/} \\
2 & smem 供数:descriptor 与 swizzle & P2.3(S042--S070) &
\texttt{cuda/m2\_smem/} \\
3 & sm100:tcgen05 & P2.4(S071--S093) & \texttt{cuda/m3\_tcgen05/} \\
4 & 完整 GEMM 四步走 & S084 + Session 4 & \texttt{cuda/m4\_gemm/} \\
5 & 低精度与 block scaling & P3(S094--S111) &
\texttt{cuda/m5\_lowprec/}、\texttt{kernels/} \\
6 & TileLang 对照 & 收尾(S112) & (复用 assignment01) \\
Team & FlashKDA / MSA decode & --- & \texttt{team/} \\
\end{longtable}

\section{环境与峰值}\label{ux73afux5883ux4e0eux5cf0ux503c}

先确认工具链和卡都能发出 tensor core 指令：编译一个最小tensor
core程序，测试卡的峰值性能，后面就可以用它来计算各个实现的性能达成率。

\begin{reading}

Session3课件S008-S021；PTX ISA 的 mma 指令一节(形状与 dtype 表)；你手里
每块卡的官方 datasheet(或 whitepaper)。

\end{reading}

\prob{0.1}{HANDS-ON}[cuda/m0\_env/01\_first\_mma.cu]

编译并运行最小 Tensor Core 程序，它使用一个 warp 执行一条 m16n8k16 的
mma.sync 指令，并与 CPU 计算结果进行比较。

\begin{verbatim}
cd assignment02/cuda
make run/m0_env/01_first_mma
\end{verbatim}

在你能使用的 GPU 上分别运行该程序（5090 使用 ARCH=120a make \ldots，B300
使用默认配置）。尝试使用不匹配的 ARCH 编译运行一次，记录现象，并结合
assignment01 Module 8 中 fatbin/JIT 的内容解释原因。可使用 make
ptx/m0\_env/01\_first\_mma 查看生成的 PTX。

\begin{answer}

硬件：RTX 4060 Laptop GPU，compute capability 8.9。\texttt{nvcc} 12.9。

默认 \texttt{make\ run/m0\_env/01\_first\_mma} 使用
\texttt{ARCH=100f}，即
\texttt{-gencode\ arch=compute\_100f,code=sm\_100f}。编译成功，运行在
\texttt{CUDA\_CHECK\_KERNEL()}（\texttt{01\_first\_mma.cu:76}）处失败：

\begin{verbatim}
CUDA error cudaErrorNoKernelImageForDevice at m0_env/01_first_mma.cu:76:
no kernel image is available for execution on the device
\end{verbatim}

\texttt{cudaMalloc} / \texttt{cudaMemcpy} 不需要 device
image，所以能走到 launch 之后
才报错。\texttt{make\ ptx/m0\_env/01\_first\_mma} 生成的 PTX 写着
\texttt{.target\ sm\_100f}， fatbin 里只有 Blackwell 的 SASS，没有 Ada
的 cubin，也没有可 JIT 到 sm\_89 的低架构 PTX。

匹配本机后通过：

\begin{verbatim}
ARCH=89 make -B run/m0_env/01_first_mma
D[0][0]=2 D[0][7]=2 D[15][0]=2 D[15][7]=2
PASS
\end{verbatim}

这与 assignment01 Module 8 的 sassonly / ptxonly 实验同一机制：driver
只从 fatbin 里挑与当前 compute capability 匹配的 SASS；对不上且没有 PTX
fallback 时，JIT 无法兜底，launch 报
\texttt{cudaErrorNoKernelImageForDevice}。 Makefile 用显式
\texttt{-gencode} 只嵌一种 \texttt{sm\_\$(ARCH)}，默认面向 B300，本机
4060 必须改 \texttt{ARCH=89}。5090 用 \texttt{ARCH=120a}，B300 用默认
\texttt{100f}。

\end{answer}

\prob{0.2}{DERIVE}

推导你所使用 GPU 的 Tensor Core 理论峰值。参考课上 A100
的推导方法（S018--S019），分别计算 5090 和 B300 的 bf16 峰值，并根据
dtype 宽度关系估算 fp8 / fp4 峰值。

开始计算前先明确采用的口径，包括 dense 或 sparse、boost 或 base
频率、FMA 是否计作 2 FLOP 等。完成推导后，再与 datasheet
中的官方数据进行对照；如果结果存在差异，需要说明差异来自哪一项口径。

\begin{longtable}[]{@{}lll@{}}
\toprule\noalign{}
量 & 5090 & B300 \\
\midrule\noalign{}
\endhead
\bottomrule\noalign{}
\endlastfoot
bf16 FLOP/cycle/SM & & \\
bf16 峰值(TFLOPS) & & \\
fp8 峰值(TFLOPS) & & \\
fp4 峰值(TFLOPS) & & \\
datasheet 对照值与口径差异 & & \\
HBM/GDDR 带宽(GB/s) & & \\
机器平衡点(FLOP/byte，bf16) & & \\
\end{longtable}

根据 bf16 峰值和显存带宽计算机器平衡点（FLOP/byte），并与单条 mma
的计算强度（S016，m16n8k16 fp16 为 3.2
FLOP/byte）比较。思考两者之间的差距意味着什么，以及为什么后续 M2--M4
需要从数据供给路径入手优化。

\prob{0.3}{CONCEPT}

判断下列说法是否正确，并给出一句理由。

\begin{enumerate}
\def\labelenumi{(\alph{enumi})}
\item
  一条 mma 的计算强度，分子是 \(2MNK\)，分母按 A、B 读入与 D 写回
  的字节总和计(S016 的口径)。
\item
  mma.sync 是 warp 级协作指令:32 个 lane 各持 fragment 的一部分， 要求全
  warp 一致地执行这条指令；有 lane 发散时行为未定义。
\item
  增大 mma 的形状 M/N/K 能提高单条指令的计算强度，而且没有代价，
  所以指令形状越大越好。
\item
  只要单条 mma 的计算强度低于机器平衡点，GEMM kernel 就不可能逼近
  计算峰值。
\end{enumerate}

\section{sm80:fragment 与 mma.sync}\label{sm80fragment-ux4e0e-mma.sync}

课上对 m16n8k16 fp16 推过 fragment 公式、手搓过单 tile mma
(C03--C07)。现在你手推一遍这个 m16n8k32 fp8 shape， 再看 ldmatrix
到底发挥了什么作用。

\begin{reading}

课件 S025--S041、C03--C08；PTX ISA 的 ``Matrix Fragments for
mma.m16n8k32'' 一节与 ``Warp-level matrix load instruction: ldmatrix''
一节；fp8 转换用 \texttt{cuda\_fp8.h}(\texttt{\_\_nv\_fp8\_e4m3})。

\end{reading}

\prob{1.1}{DERIVE}[cuda/m1\_sm80/01\_fragment\_map.cu]

根据 PTX 文档中的 fragment 布局，推导
\texttt{mma.sync.aligned.m16n8k32.row.col.f32.e4m3.e4m3.f32} 对应的 A、B
fragment 映射公式，并完成文件中的四个函数。 判测使用纯 host 的 32-lane
真值表比对，无需 GPU 即可完成：

\begin{verbatim}
cd assignment02/cuda
make run/m1_sm80/01_fragment_map
\end{verbatim}

附加问题（写入报告）： A 的同一个 b32 寄存器中的 4 个 fp8
元素沿矩阵哪个方向相邻？ 这个布局对 1.4 中使用 ldmatrix load
有什么影响？

\prob{1.2}{DEBUG}[cuda/m1\_sm80/02\_bug\_fragment.cu]

这个程序发一条 m16n8k16 fp16 mma，判测会 FAIL。先运行一遍，后改动:

\begin{enumerate}
\def\labelenumi{(\alph{enumi})}
\item
  描述症状：D 的哪些位置错、错成了什么(和对的部分是什么关系)；
\item
  修好它，并解释错的是哪个 fragment
  的哪部分映射，为什么恰好产生(a)的症状。
\end{enumerate}

\begin{verbatim}
cd assignment02/cuda
make run/m1_sm80/02_bug_fragment
\end{verbatim}

\begin{capstone}{prob 1.3(FROM-SCRATCH):手写单 tile fp8 mma}

在 \texttt{cuda/m1\_sm80/03\_mma\_fp8.cu} 中从零实现一个单 tile 的 fp8
mma，不提供代码骨架。使用 \texttt{m16n8k32} e4m3 mma、f32
累加，并手动完成 fragment 装载（本题不使用 ldmatrix），最后与 CPU
参考结果进行严格相等比较。

要求如下：

输入使用小整数，保证转换为 e4m3 后可以精确表示；

程序接受一个 seed 参数，例如 \texttt{./prog\ 123}；

输出必须以 \texttt{PASS} 或 \texttt{MISMATCH} / \texttt{FAIL} 开头；

fp8 与 float 的互转直接使用 \texttt{cuda\_fp8.h}。

1.1 中推导的 fragment
映射公式会直接用在这里。映射只要有一处错误，随机数据的判测就无法通过。

判测脚本会使用五个 seed，全部通过才算完成：

\begin{verbatim}
cd assignment02/cuda/m1_sm80
./judge_mma_fp8.sh 03_mma_fp8.cu
\end{verbatim}

\end{capstone}

\prob{1.4}{MODIFY}[cuda/m1\_sm80/04\_ldmatrix.cu]

在给定骨架中保留 1.3 的手工装载方式，并另外实现一条使用
\texttt{ldmatrix} 的装载路径。两种实现需要共存，并分别通过判测。

根据 PTX 文档选择合适的 \texttt{ldmatrix}
变体（\texttt{.x1/.x2/.x4}，以及是否使用 \texttt{.trans}）。B 矩阵在
shared memory 中的布局还需要满足 \texttt{ldmatrix} 的 16 byte
行地址要求，具体约束见文件头说明。

两条路径都通过判测后，使用 \texttt{make\ ptx/m1\_sm80/04\_ldmatrix} 或
\texttt{nvdisasm} 查看生成的指令，分别统计 \texttt{smem\ →\ fragment}
阶段的装载指令和地址计算指令数量，并回答：

\begin{enumerate}
\def\labelenumi{(\alph{enumi})}
\item
  \texttt{ldmatrix} 省掉了手工装载中的哪些工作？
\item
  为什么这些工作在手工装载路径中无法避免？
\end{enumerate}

\prob{1.5}{EXPERIMENT}[cuda/m1\_sm80/05\_ldsm\_stride.cu]

观察行跨度对 \texttt{ldmatrix} 的影响。对同一个 16×16 fp16
tile，分别使用 32 B、64 B、128 B 和 128+16 B（padding）四种行跨度。

先根据课上介绍的 bank 模型，预测四种情况下执行一次 \texttt{ldmatrix}
所需的 wavefront 数及其比例，再运行程序并使用 Nsight Compute 测量：

\begin{Shaded}
\begin{Highlighting}[]
\BuiltInTok{cd}\NormalTok{ assignment02/cuda}
\FunctionTok{make}\NormalTok{ run/m1\_sm80/05\_ldsm\_stride}
\ExtensionTok{ncu} \AttributeTok{{-}{-}metrics}\NormalTok{ l1tex\_\_data\_pipe\_lsu\_wavefronts\_mem\_shared\_op\_ld.sum,l1tex\_\_data\_bank\_conflicts\_pipe\_lsu\_mem\_shared\_op\_ld.sum ./bin/m1\_sm80/05\_ldsm\_stride}
\end{Highlighting}
\end{Shaded}

\begin{longtable}[]{@{}
  >{\raggedright\arraybackslash}p{(\linewidth - 8\tabcolsep) * \real{0.2000}}
  >{\raggedright\arraybackslash}p{(\linewidth - 8\tabcolsep) * \real{0.2000}}
  >{\raggedright\arraybackslash}p{(\linewidth - 8\tabcolsep) * \real{0.2000}}
  >{\raggedright\arraybackslash}p{(\linewidth - 8\tabcolsep) * \real{0.2000}}
  >{\raggedright\arraybackslash}p{(\linewidth - 8\tabcolsep) * \real{0.2000}}@{}}
\toprule\noalign{}
\begin{minipage}[b]{\linewidth}\raggedright
档位
\end{minipage} & \begin{minipage}[b]{\linewidth}\raggedright
预测 wavefront 比
\end{minipage} & \begin{minipage}[b]{\linewidth}\raggedright
实测 wavefront
\end{minipage} & \begin{minipage}[b]{\linewidth}\raggedright
实测 conflict
\end{minipage} & \begin{minipage}[b]{\linewidth}\raggedright
平均 cycle
\end{minipage} \\
\midrule\noalign{}
\endhead
\bottomrule\noalign{}
\endlastfoot
32 B & & & & \\
64 B & & & & \\
128 B & & & & \\
128+16 B & & & & \\
\end{longtable}

比较预测与实测结果：哪一种行跨度使 wavefront 数增加到 4 倍？ wavefront
的比例应与 bank 模型一致，但实际耗时的差距通常没有这么大。 结合 8 个
warp 的占用情况，解释为什么 wavefront 增加 4 倍并不会使总耗时也增加 4
倍。

\begin{lookback}

1.3--1.5 关注的是同一个问题：怎样把数据从 shared memory 装入 Tensor Core
的 fragment。1.3 手动完成每个 lane 的装载，1.4 使用 \texttt{ldmatrix}
简化这一步，而 1.5 说明即使使用 \texttt{ldmatrix}，shared memory 的 bank
conflict 仍然会影响实际效率。

后面的模块会继续沿着这条数据供给路径展开：M2 中使用 descriptor
描述布局，M4 中进一步使用 TMA 和 pipeline 完成数据搬运。

\end{lookback}

\section{smem 供数:descriptor 与
swizzle}\label{smem-ux4f9bux6570descriptor-ux4e0e-swizzle}

sm90 起，tensor core 从 smem 取数不再经过 fragment 装载指令，而是 读一个
64 位 descriptor 按 canonical 布局自取。descriptor 与 swizzle 的格式在
sm100 的 tcgen05 上原样沿用，本模块推的每个公式都是 M3、 M4 的直接组件。

\begin{reading}

课件 S042--S070(proxy:S051；core matrix 与 LBO/SBO:S057--S061；
swizzle:S062--S065)；PTX ISA 的 ``Asynchronous Warpgroup MMA Shared
Memory Layout'' 与 swizzling 小节。

\end{reading}

\prob{2.1}{CONCEPT}

\begin{enumerate}
\def\labelenumi{(\alph{enumi})}
\tightlist
\item
  一个 warpgroup 使用 wgmma 读取刚写入 shared memory
  的数据。将下面六个操作排成正确顺序，并说明每一步用于避免哪两个参与者之间的哪种乱序：
\end{enumerate}

\texttt{wgmma.mma\_async} / \texttt{st.shared} /
\texttt{wgmma.commit\_group} / \texttt{fence.proxy.async} /
\texttt{wgmma.fence} / \texttt{wgmma.wait\_group}

\begin{enumerate}
\def\labelenumi{(\alph{enumi})}
\setcounter{enumi}{1}
\tightlist
\item
  判断下列说法是否正确，并给出一句理由。
\end{enumerate}

\begin{enumerate}
\def\labelenumi{\arabic{enumi}.}
\tightlist
\item
  \texttt{fence.proxy.async} 是 wgmma 专属的指令，TMA 与 tcgen05
  的场景不需要它。
\item
  \texttt{wgmma.commit\_group} 会阻塞，直到它之前发射的 wgmma 全部完成。
\item
  不加 \texttt{fence.proxy.async} 时，wgmma 可能读到 shared memory
  中的旧值，因为 \texttt{st.shared} 的写经过 generic proxy，而 wgmma
  的读经过 async proxy。
\end{enumerate}

\prob{2.2}{DERIVE}[cuda/m2\_smem/02\_descriptor.cu]

根据文件头给出的位域，实现 SM100 的 64 位 smem matrix descriptor
编码函数，并分别对下面三种情况推导 LBO、SBO 和 layout：

\begin{itemize}
\tightlist
\item
  K-major，无 swizzle；
\item
  K-major，128B swizzle；
\item
  MN-major，128B swizzle。
\end{itemize}

判测为纯 host，无需 GPU。测试使用的描述符真值已经在 B300 上通过实际
tcgen05 指令验证，3.2 中也会使用这三组描述符：

\begin{verbatim}
cd assignment02/cuda
make run/m2_smem/02_descriptor
\end{verbatim}

场景 2 和场景 3 最终得到的 descriptor 相同。在报告中回答：MN-major 与
K-major 的区别体现在哪里？

\prob{2.3}{FROM-SCRATCH}[cuda/m2\_smem/03\_swizzle.cu]

实现 128B、64B 和 32B 三种 swizzle 模式的地址映射函数，将逻辑坐标映射到
atom 内的物理字节偏移。

课件 S064 已经推导了 128B swizzle 的地址位异或关系；64B 和 32B
两种模式需要根据 PTX ISA 的 swizzling 一节自行推导。

判测分为两部分：首先检查映射是否为双射，然后检查列访问是否存在 bank
conflict。单纯使用 padding
或恒等映射虽然可能通过第一项，但无法通过第二项。判测同样为纯 host：

\begin{verbatim}
cd assignment02/cuda
make run/m2_smem/03_swizzle
\end{verbatim}

本题实现的 \texttt{swizzle\_128B} 会在 3.2 中直接用于 shared memory
staging，并接受实际硬件上的 GEMM 判测。若 3.2 或 4.1
出现问题，可以先运行本题的判测，排除 swizzle 布局错误。

\begin{lookback}

课件 C12--C14 给出了使用 wgmma 实现单 tile
的完整示例，本作业不单独设置对应题目。wgmma 是 sm\_90a
专属指令，5090（sm120）和
B300（sm100）均无法运行；有兴趣的同学可以在集群 H 卡上自行尝试。

本模块重点保留 wgmma 路径中可以继续使用的两个部分：descriptor 与
swizzle。2.2、2.3 先分别完成它们的推导与判测，后面的 M3
会在实际硬件上继续使用。

\end{lookback}

\section{sm100：tcgen05}\label{sm100tcgen05}

在 sm100 上，tcgen05 将累加器放入 TMEM，并使用 mbarrier
处理异步完成通知。本模块会先在 B300 上完成一个单 tile
GEMM，再通过后续实验理解 TMEM、mbarrier 和 2-CTA MMA 的使用方式。

\begin{reading}

课件 S071--S093（TMEM：S073--S075；mma 与 idesc：S076--S079；
mbarrier：S080--S084；ld 与同步：S085--S086；七步流程：S087/F27；
2-CTA：S091）；PTX ISA 的 tcgen05 一族小节（alloc / mma / commit / ld /
fence）与 mbarrier 一节。

\end{reading}

\prob{3.1}{CONCEPT}

判断下列说法是否正确，并给出一句理由。其中 (d) 需要写出计算过程。

\begin{enumerate}
\def\labelenumi{(\alph{enumi})}
\item
  \texttt{tcgen05.ld} 读取 TMEM 时，每个 warp 只能读取自己对应的 32 条
  lane，warp 之间不能互相读取。
\item
  与 \texttt{mma.sync} 由 warp 协作、wgmma 由 warpgroup
  协作不同，\texttt{tcgen05.mma} 由单个线程发射，随后由硬件异步执行。
\item
  TMEM 中的累加结果可以直接通过 TMA 搬回 global
  memory，不需要经过寄存器。
\item
  TMEM 每个 SM 包含 128 lane × 512 column × 4 B；一个 m128n256 的 f32
  accumulator 恰好占用其中一半。
\item
  \texttt{tcgen05.commit} 会阻塞直到之前发射的 mma 全部完成，因此 commit
  返回后即可安全读取 TMEM。
\end{enumerate}

\begin{capstone}{prob 3.2(FROM-SCRATCH):tcgen05 单 tile GEMM \hfill {\footnotesize\ttfamily cuda/m3\_tcgen05/02\_single\_tile.cu}}

从零实现一个 tcgen05 单 tile GEMM：计算 m128n64k64 的 bf16 矩阵乘，使用
f32 累加、\texttt{cta\_group::1}，并由一个 128 线程的 block 完成。

数据按照下面的路径流动：

\texttt{global\ →\ smem（K-major\ +\ 128B\ swizzle）→\ tcgen05.mma\ →\ TMEM\ →\ tcgen05.ld\ →\ global}

代码骨架只保留课件 F27 中的七步流程注释。descriptor 使用 2.2
中实现的编码方式，shared memory staging 使用 2.3 中的
\texttt{swizzle\_128B}。TMEM alloc、mbarrier、idesc 以及 mma / ld
的具体写法需要根据 PTX ISA 完成，相关语义可参考课件 C15--C21。

\texttt{fence.proxy.async} 的位置以及 \texttt{tcgen05.ld} 的 warp
可见范围已经在文件头中给出。

判测：

\begin{verbatim}
cd assignment02/cuda
make run/m3_tcgen05/02_single_tile
cd m3_tcgen05 && ./judge_tile.sh
\end{verbatim}

在正确版本通过后，故意去掉 \texttt{fence.proxy.async}
再运行一次，记录出现的现象并写入报告。结合 2.1(a) 中的排序问题解释原因。

\end{capstone}

\prob{3.3}{DEBUG}[cuda/m3\_tcgen05/03\_bug\_mbarrier.cu]

这个程序连续发射多轮 mma，并在每一轮后读取结果，用来模拟 M4 中沿 K
维流式计算的过程。当前程序的 mbarrier
使用存在错误，判测带有超时机制，因此程序挂死本身也是需要观察的现象。

先运行：

\begin{verbatim}
cd assignment02/cuda
make bin/m3_tcgen05/03_bug_mbarrier
cd m3_tcgen05 && ./judge_mbar.sh
\end{verbatim}

\begin{enumerate}
\def\labelenumi{(\alph{enumi})}
\item
  分别记录 \texttt{rounds\ =\ 1}、\texttt{2}、\texttt{4}
  时的运行结果以及问题出现的条件。
\item
  修复程序，并画出修改前后 mbarrier 的状态变化，包括 phase 和 arrival
  count。指出错误版本在哪一次等待中过早或过晚放行，并解释为什么会产生
  (a) 中观察到的现象。
\end{enumerate}

提示：注意 \texttt{tcgen05.ld} 与尚未完成的 mma 之间的关系。

\prob{3.4}{EXPERIMENT}[cuda/m3\_tcgen05/04\_cta\_pair.cu]

比较 \texttt{cta\_group::1} 和 \texttt{cta\_group::2} 完成同一个
m256n64k64 任务时的数据开销。

\texttt{cta\_group::1} 使用两个独立 block；\texttt{cta\_group::2}
使用一个 cluster 发射一条 M=256 的 mma，两个 CTA 分别保存 B
矩阵的一部分。

先预测，再运行实验：

\begin{enumerate}
\def\labelenumi{(\alph{enumi})}
\item
  在 \texttt{cta\_group::2} 中，每个 CTA 所需的 B shared memory 是
  \texttt{cta\_group::1} 的多少？TMEM 的占用又如何变化？程序会打印
  shared memory 用量，用它验证你的预测。
\item
  使用 Nsight Compute 比较两种实现的 shared memory 总流量。
\item
  \texttt{cta\_group::2} 节省下来的 shared memory 容量，在 M4 的
  pipeline 中可以用来做什么？
\item
  \texttt{cta\_group::2} 依赖 sm90 引入的哪一种硬件机制？5090 不支持
  2-CTA MMA，结合 0.2
  中整理的硬件参数，说明为什么这类机制更常出现在数据中心 GPU 上。
\end{enumerate}

运行：

\begin{verbatim}
cd assignment02/cuda
make run/m3_tcgen05/04_cta_pair
\end{verbatim}

在这个单 tile
实验中，两种实现的运行时间差异处于噪声范围内，因此不要用耗时判断优劣。主要比较
shared memory 用量以及 Nsight Compute 测得的流量。

\section{完整 GEMM}\label{ux5b8cux6574-gemm}

本模块将 3.2 的单 tile 实现扩展为完整的 B300 GEMM，并依次加入
tiling、TMA 和多级 pipeline。实验统一使用 4096³ 的 bf16 GEMM， tile
大小固定为 128×64×64，并与 cuBLAS 结果进行严格相等比较。

\begin{reading}

课件 S084（pipeline 骨架）、S087；Session 4 讲义对应章节；PTX ISA 的
\texttt{cp.async.bulk.tensor}、mbarrier（\texttt{expect\_tx}）小节；CUDA
Driver API 的 \texttt{cuTensorMapEncodeTiled}。

\end{reading}

下面的表格贯穿 4.1--4.3。第 0 行需要在同一块 GPU 上重新运行 assignment01
Bonus 中的 naive matmul。由于该实现使用 fp32，只比较性能量级。

\begin{longtable}[]{@{}llll@{}}
\toprule\noalign{}
实现 & TFLOPS & 对 cuBLAS 达成率 & 一句话：时间主要花在哪 \\
\midrule\noalign{}
\endhead
\bottomrule\noalign{}
\endlastfoot
naive（assignment01，fp32） & & & \\
4.1 tiled & & & \\
4.2 TMA & & & \\
4.3 pipeline（S=3） & & & \\
cuBLAS & & 100\% & \\
\end{longtable}

\prob{4.1}{FROM-SCRATCH}[cuda/m4\_gemm/01\_tiled.cu]

将 3.2 的单 tile 实现扩展为完整 GEMM：使用 grid 覆盖所有输出 tile， 并沿
K 维循环完成累加。数据 staging 仍然使用 \texttt{st.shared} 和 swizzle。

相对 3.2，本题主要新增 tile 映射以及 K 循环中的同步和累加。
具体步骤见文件头。

\begin{verbatim}
cd assignment02/cuda
make run/m4_gemm/01_tiled
\end{verbatim}

通过判测后，填写性能表中 \texttt{4.1\ tiled} 一行。结合 0.2 中计算的机器
平衡点，判断此时性能主要受哪一环节限制。

\prob{4.2}{MODIFY}[cuda/m4\_gemm/02\_tma.cu]

将 4.1 中的 shared memory staging 改为 TMA。

host 端使用 \texttt{cuTensorMapEncodeTiled} 创建 tensor map；kernel
中使用 \texttt{cp.async.bulk.tensor} 搭配 mbarrier 的
\texttt{expect\_tx} 完成搬运。 本题仍然使用单缓冲。

tensor map 的参数以及同步方式的变化已经在文件头中给出。 swizzle 由 TMA
根据 tensor map 自动完成，原有 descriptor 不需要修改字段。

\begin{verbatim}
cd assignment02/cuda
make run/m4_gemm/02_tma
\end{verbatim}

通过判测后，填写性能表中 \texttt{4.2\ TMA} 一行。

比较 4.1 与 4.2，并回答：4.1 中 shared memory staging 的开销由哪些
部分组成？改用 TMA 后，其中哪些工作不再由普通 CUDA 指令完成？

使用 Nsight Compute 辅助分析，观察 4.1 中 SM 时间主要消耗在哪些部分。

\begin{capstone}{prob 4.3(FROM-SCRATCH):多级流水 \hfill {\footnotesize\ttfamily cuda/m4\_gemm/03\_pipeline.cu}}

将 4.2 的单缓冲改为 \texttt{STAGES} 级循环缓冲，使后续 K tile 的 TMA
预取能够与当前 tile 的 mma 计算重叠。

\texttt{STAGES} 为编译参数，例如：

\begin{verbatim}
STAGES=4 make -B run/m4_gemm/03_pipeline
\end{verbatim}

这里的 \texttt{-B} 不能省略，否则修改 \texttt{STAGES}
后可能不会重新编译。

文件头给出了每个 stage 使用双 mbarrier 的基本结构，同时包含一个需要
处理的 pipeline hazard：强制发射与机会式预取之间的边界处理错误会导致
死锁。在该错误下，1024³ 可能偶尔通过，而 4096³ 会稳定暴露问题。
开始实现前先阅读文件头说明。

本题不要求 warp specialization、persistent kernel 或 epilogue 融合，
也不设置 cuBLAS 达成率门槛。重点是流水实现是否正确，以及能否根据实验
结果解释性能变化。

\begin{verbatim}
cd assignment02/cuda
make run/m4_gemm/03_pipeline
cd m4_gemm && ./sweep_stages.sh
\end{verbatim}

完成以下内容：

\begin{enumerate}
\def\labelenumi{\arabic{enumi}.}
\item
  使用 \texttt{S=3} 的结果填写性能表中 \texttt{4.3\ pipeline} 一行。
\item
  完成不同 stage 数的性能测试：

  \begin{longtable}[]{@{}lllll@{}}
  \toprule\noalign{}
  形状 & S=2 & S=3 & S=4 & S=6 \\
  \midrule\noalign{}
  \endhead
  \bottomrule\noalign{}
  \endlastfoot
  4096³ & & & & \\
  256 × 4096 × 16384 & & & & \\
  \end{longtable}

  比较两个形状对 \texttt{STAGES} 的敏感程度，并结合 shared memory 用量、
  每个 SM 可同时驻留的 block 数以及 block 间并发能够隐藏的延迟进行解释。
\item
  任选一个 \texttt{STAGES}，画出稳态阶段各 stage 中 TMA 与 mma
  的流水时空图。
\item
  回答：

  \begin{enumerate}
  \def\labelenumii{(\alph{enumii})}
  \item
    从 4.1 到 4.3，主要瓶颈发生了哪些变化？
  \item
    梯子表中每一级优化分别减少了哪部分开销？
  \item
    如果继续增大 tile 或增加 stage 数，shared memory 与 TMEM
    哪一个会先成为容量限制？结合 3.4(c) 的结果说明。
  \end{enumerate}
\end{enumerate}

\end{capstone}

\prob[Optional]{4.4}{FROM-SCRATCH}

任选一个方向继续优化：

\begin{enumerate}
\def\labelenumi{(\alph{enumi})}
\item
  实现 4.3 的 \texttt{cta\_group::2} 版本，利用 B shared memory
  用量的减少 增加 pipeline 深度或扩大 tile；
\item
  在 5090 上使用 \texttt{mma.sync} 实现相同的 tiling，并与 B300
  的结果比较；
\item
  自由优化当前 kernel，提高相对 cuBLAS 的性能，并记录每一步优化
  解决了什么问题。
\end{enumerate}

\prob{4.5}{EXPERIMENT}[cuda/m4\_gemm/05\_thin\_gemm.cu]

观察矩阵形状较窄时 Tensor Core GEMM 的性能变化。

实验形状取自 vLLM 主树中 Kimi K3 的 decode GEMM dispatch 表：

\texttt{vllm/models/kimi\_k3/nvidia/low\_latency\_gemm.py}

使用其中七个投影层对应的 N、K，并对 M 取：

\(\{1, 8, 16, 64, 256, 1024, 4096, 16384, 65536\}\)

其中 N 和 K 由权重形状决定，变化的 M 表示一次 GEMM 处理的 token 数量。
较小的 M 对应 decode batch；较大的 M 可以对应 chunked prefill 中单次
处理的 chunk。vLLM 在 \(M \le 16\) 时会放弃 cuBLAS / Tensor Core 路径，
改用 CUDA Core FMA 的 skinny kernel。

运行程序前，先对每个形状计算 arithmetic intensity：

\[
AI = \frac{2MNK}{2MK + 2NK + 2MN}
\]

将结果与 0.2 中得到的机器平衡点比较，并计算对应的理论性能上限：

\begin{itemize}
\tightlist
\item
  compute bound：Tensor Core 峰值；
\item
  memory bound：\(AI \times\) 显存带宽。
\end{itemize}

然后运行：

\begin{verbatim}
cd assignment02/cuda
make bin/m4_gemm/05_thin_gemm
./bin/m4_gemm/05_thin_gemm <峰值TFLOPS> <带宽GB/s>
\end{verbatim}

其中两个参数使用 0.2 中得到的数值。

程序会输出 TFLOPS、GB/s、AI 以及相对于 compute roof 和 memory roof
的达成率。根据结果回答：

\begin{enumerate}
\def\labelenumi{(\alph{enumi})}
\item
  描述性能随 M 变化的趋势。M 较小时，Tensor Core 性能从什么时候
  开始明显下降？M
  增大到什么范围后进入相对稳定的平台？平台上的达成率是多少？
\item
  对性能下降的形状，比较 compute roof 和 memory roof 两种达成率。
  哪些形状主要受到显存带宽限制？
\item
  \texttt{f\_b\_proj}（K=128）中两个 roof
  的达成率都较低。结合矩阵形状分析， 它的限制来自哪里？
\item
  根据上述结果解释 vLLM 在 \(M \le 16\) 时选择 skinny CUDA Core kernel
  的原因。
\end{enumerate}

\texttt{in\_proj\_qkvgfab} 对应 KDA 的输入投影。完成团队题 C1/C2
的同学可以 使用这一行的实验结果作为后续分析的参考。

\begin{lookback}

4.1--4.3 从同一个 GEMM 出发，依次加入 tiling、TMA 和多级 pipeline。
回顾梯子表时，重点不是最终的 cuBLAS 达成率，而是能够说明每一步优化
减少了什么开销，以及新的瓶颈出现在哪里。

assignment01 Bonus 中 naive matmul 与 cuBLAS 之间的部分差距，现在已经
可以从 Tensor Core、异步数据搬运和软件流水三个方面解释。进一步的 warp
specialization、更大的 tile、epilogue 融合和 persistent kernel
等优化不在本作业要求范围内。

4.5 则说明 Tensor Core 的收益同样依赖矩阵形状。当 M 很小时，即使 kernel
本身使用了前面的优化方法，也不一定能够有效利用 Tensor Core。

\end{lookback}

\section{低精度与 block
scaling}\label{ux4f4eux7cbeux5ea6ux4e0e-block-scaling}

降低数值精度可以换取更高的 Tensor Core
吞吐，但代价是可表示的动态范围变小。本模块先从 per-tensor scale 遇到
outlier 时的问题入手，再分析 block scaling 的代数约束，最后完成一条
NVFP4 量化通路，并通过融合实验观察低精度 kernel 的实际性能。

\begin{reading}

课件 S094--S111（fp8 格式：S096；per-tensor 与 outlier：S097--S099；
block scaling 约束：S100--S101；fp4 与 NVFP4：S106--S108；SF 布局：
S108）；\texttt{cuda\_fp4.h} 中的 \texttt{\_\_nv\_fp4x2\_e2m1}；cuBLASLt
的 block-scaled FP4 matmul。格式约定见题面材料
\texttt{cuda/m5\_lowprec/nvfp4\_common.h}。

\end{reading}

\prob{5.1}{EXPERIMENT}[kernels/quant\_outlier.py]

观察 outlier 对 per-tensor scale 的影响。输入包含一万个 \([-1,1]\)
均匀分布的元素，并额外加入一个值为 3000 的 outlier。补全脚本中的两个
TODO，完成 per-tensor E4M3 量化与反量化：

\begin{verbatim}
cd assignment02
uv run python kernels/quant_outlier.py
\end{verbatim}

\begin{longtable}[]{@{}llllll@{}}
\toprule\noalign{}
采样点 \(x\approx\) & 0.5 & 0.1 & 0.01 & 0.005 & 3000 \\
\midrule\noalign{}
\endhead
\bottomrule\noalign{}
\endlastfoot
相对误差 & & & & & \\
\end{longtable}

根据实验结果回答：

\begin{enumerate}
\def\labelenumi{(\alph{enumi})}
\item
  去掉 outlier 后重新量化，\(x\approx0.5\) 处的误差变化多少倍？
\item
  找出输入被量化为 0 的阈值，并写出该阈值与 scale 的关系式。
\item
  改用 1×128 的 per-block scale 后，包含 outlier 的 block 与不包含
  outlier 的 block 分别有什么变化？
\end{enumerate}

\prob{5.2}{DERIVE}[kernels/block\_scale\_sim.py]

block scaling 的代数关键是沿着哪个方向分段 scale 乘积在 K
归约中才能保持常数。补全两个 fp64 模拟函数:

\begin{itemize}
\tightlist
\item
  \texttt{gemm\_scale\_per\_row\_col}:A 每行一个 scale，B 每个输出列一个
  scale。scale 乘积在整个点积中不变，完整归约后只用乘回一次。
\item
  \texttt{gemm\_scale\_along\_k}:A 和 B 的 scale 每 128 个 K 元素改变
  一次。每个 K block 的 partial sum 要分别乘回该段的 scale 乘积，
  然后进行累加。
\end{itemize}

文件还提供了错误范例
\texttt{gemm\_scale\_along\_k\_one\_restore}：它先把 所有归一化 partial
sum 相加，最后只乘回第一段的 scale。

判测:

\begin{verbatim}
cd assignment02
uv run pytest tests/test_block_scale.py
\end{verbatim}

两个正确函数只要在 fp64 容差内与直接 GEMM 等价，不要要求
bit-exact（因为分段会改变浮点加法的分组）。然后回答:

\begin{enumerate}
\def\labelenumi{(\alph{enumi})}
\item
  用两行代数式分别写出 row/column scale 为什么可以在整个 点积外乘回，而
  K-block scale 为什么必须逐段乘回。
\item
  \href{https://docs.nvidia.com/cutlass/latest/media/docs/cpp/blackwell_functionality.html}{NVIDIA
  CUTLASS 的 Blackwell SM100 GEMM 说明} 把 A 的 scale 布局写成
  \(M \times \lceil K/SV \rceil\)，B 写成
  \(N \times \lceil K/SV \rceil\)，每个 scale 负责连续的 16 或 32 个 K
  元素。结合 GEMM 内循环、连续供数和 Tensor Core 指令语义，说明硬件
  为什么把 block scale 与 K 归约对齐。
\item
  DeepSeek-V3 的 weight 128×128 表示 scale 还在输出通道方向上
  共享，但每个组仍覆盖一段 K；NVFP4 则每 16 个 K 元素一组。 结合 5.1(c)
  的误差，说明粒度 16 相对粒度 128 有什么优势，又增加了 多少 scale
  metadata 与供数复杂度。
\end{enumerate}

\prob{5.3}{FROM-SCRATCH}[cuda/m5\_lowprec/]

完成一条 NVFP4 量化通路。开始前先阅读 \texttt{nvfp4\_common.h}
中给出的格式约定：

\begin{itemize}
\tightlist
\item
  每个量化组包含 K 方向连续的 16 个元素；
\item
  SF 使用 \texttt{amax\ /\ 6}，并以 e4m3 保存；
\item
  SF 按 Tensor Core 消费要求的 swizzled 布局存放，字节偏移公式已经给出。
\end{itemize}

\paragraph{(a) E2M1 编码}\label{a-e2m1-ux7f16ux7801}

在 \texttt{e2m1\_encode.h} 中实现 host/device 通用的 E2M1
round-to-nearest-even 编码器。

判测会在 GPU 上使用 \texttt{cuda\_fp4.h} 中的
\texttt{\_\_nv\_fp4x2\_e2m1}
转换同一批候选值，包括所有中点、边界以及采样值，并与自己的编码结果
逐位比较：

\begin{verbatim}
cd assignment02/cuda
make run/m5_lowprec/03a_encode_check
\end{verbatim}

\paragraph{(b) NVFP4 quant kernel}\label{b-nvfp4-quant-kernel}

在 \texttt{nvfp4\_quant\_kernel.h} 中实现 quant kernel，完成

\texttt{bf16\ →\ e2m1\ 打包\ +\ e4m3\ SF\ +\ swizzled\ SF\ 布局}

设备侧的 E2M1 转换直接使用 \texttt{\_\_nv\_fp4x2\_e2m1}。

第一层判测要求输出逐 byte 与 host 参考结果相等；host 参考使用你在 5.3(a)
中实现的 E2M1 编码器生成。随后将量化结果原样交给 cuBLASLt 的 FP4
matmul，检查生成的数据和 SF 布局能否被实际消费：

\begin{verbatim}
make run/m5_lowprec/03b_nvfp4_quant
make run/m5_lowprec/test_fp4_gemm
\end{verbatim}

正确实现下，FP4 matmul 的 \texttt{maxrel} 约为 \texttt{4e-3}，主要来自
bf16 输出的舍入误差。如果 SF 布局错位或 scale 对应到了错误的量化组，
通常会出现成块的明显数值错误，而不是仅有小幅舍入误差。

\paragraph{(c) Ceiling probe}\label{c-ceiling-probe}

在 \texttt{03c\_ceiling\_probe.cu} 中实现一个 ceiling probe。它与 quant
kernel 使用相同的访存模式，但不执行量化计算，只读取数据并通过简单 xor
后写回。

报告：

\begin{itemize}
\tightlist
\item
  ceiling probe 的 GB/s；
\item
  quant kernel 的 GB/s；
\item
  两者的比值。
\end{itemize}

结合 Nsight Compute 判断 quant kernel 距离自己的访存上限还有多远，
以及剩余差距主要来自访存还是计算。

\paragraph{(d) Optional}\label{d-optional}

使用 tcgen05 \texttt{kind::mxf4nvf4} 直接消费自己生成的 NVFP4 数据， SF
通过 TMEM 提供，并与 cuBLASLt 的结果对拍。

\begin{lookback}

5.3(c) 用来区分量化 kernel 中的访存成本和数值转换成本。作为参考，
软件实现 E2M1 RN-even 时，每个 kernel 约执行 18129 条 FSETP；
使用硬件转换后约为 1272 条 \texttt{F2FP.E2M1}，同一 kernel 的时间从 71.7
µs 降到 32.8 µs。Nsight Compute 中的瓶颈也从 SM 利用率 84.7\% 的
compute-bound 转为 memory-bound。

这组结果说明，当低精度转换由专用硬件完成后，量化 kernel 可以更接近
纯数据搬运的性能特征；5.3(c) 的 ceiling probe 就用于衡量这一差距。

Hopper 上的 fp8 累加需要软件 promotion（S102--S105，DeepGEMM 使用
双层累加循环），而 Blackwell 进一步在硬件中支持 block scaling
（S105）。本作业不单独设置对应题目，但 5.2 中的 scale 分组约束是
理解两者的共同基础。

\end{lookback}

\begin{capstone}{prob 5.4(FROM-SCRATCH):融合 rms\_norm + NVFP4}

实现融合的 \texttt{rms\_norm\ +\ NVFP4} kernel：

\[
y = \mathrm{rms\_norm}(x)\cdot w
\]

要求将结果直接量化为 NVFP4，不写回 bf16 中间结果。

本题以 vLLM fused kernel 清单中的 issue \#25179，以及两个相关实现 PR
\#36413、\#32957 为背景。相关实现的正确性没有问题，也确实将 kernel
数量从 2 个减少到了 1 个，但当时观察到的端到端收益仍处于 ±1\% 左右的
测量噪声内。本题要通过逐形状实验解释这部分收益去了哪里。

按理论字节量计算，两步实现约为 6.56 B/elem，融合后约为 2.56
B/elem，因此单纯从访存量估计可以得到约 2.56× 的理想加速。
这个数字只是上限，实际结果需要通过实验验证。

程序已经提供两步基线、正确性检查和逐形状计时框架。测试覆盖：

\begin{itemize}
\tightlist
\item
  \(M\) 从 1 到 16384；
\item
  \(K \in \{4096, 7168, 8192\}\)；
\item
  共十个测试形状。
\end{itemize}

需要完成三部分：

\begin{enumerate}
\def\labelenumi{\arabic{enumi}.}
\tightlist
\item
  实现融合 kernel，具体线程组织和 tile 结构不限；
\item
  分别调优融合实现和两步基线，使两边都使用各自合理的配置后再比较；
\item
  完成逐形状性能分析。
\end{enumerate}

第二点必须认真处理。让两步基线处于明显不利的配置下得到的加速比没有
比较意义，这也是相关上游 PR 的实验过程中暴露出的一个问题。

运行：

\begin{verbatim}
cd assignment02/cuda
make run/m5_lowprec/04_fused_rms_nvfp4
\end{verbatim}

报告逐形状结果，包括：

\begin{itemize}
\tightlist
\item
  实测加速比；
\item
  融合 kernel 相对 5.3(c) ceiling probe 的性能比例。
\end{itemize}

重点解释实测结果与 2.56× 理论上限之间的差距：不同 M 范围内主要受到
什么因素限制，并给出 Nsight Compute 数据或计算结果作为依据。

本题不设置性能门槛。实现首先需要保证正确，评分重点放在实验设计和性能归因。

Optional：根据分析得到的主要瓶颈进行一次针对性优化，重新测试并更新表格。

\end{capstone}

\prob{5.5}{CONCEPT}

比较 W4A16 + Marlin kernel（权重 int4、计算 fp16）与 5.3 中的 NVFP4
GEMM。

根据课件 S109 的分类回答：

\begin{enumerate}
\def\labelenumi{(\alph{enumi})}
\item
  两者分别属于存储量化还是计算量化？
\item
  两种方法分别节省哪些资源：显存容量、显存带宽还是计算吞吐？
\item
  在 4.5 中的小 batch decode 场景下，哪一类量化的收益更加直接？
\end{enumerate}

每问用两到三句话回答。

\section{TileLang 对照}\label{tilelang-ux5bf9ux7167}

前面的模块已经手动处理过 Tensor Core 指令、数据布局和数据搬运。
本模块通过 TileLang 的 lowering 结果观察其中哪些工作可以交给 DSL 完成。

\begin{reading}

课件 S112；assignment01 Module 7（7.5 的表、7.6 的
\texttt{kernels/tilelang\_matmul.py}）；TileLang 文档中的 lowering/debug
部分。

\end{reading}

\prob{6.1}{EXPERIMENT}

使用 assignment01 的 \texttt{kernels/tilelang\_matmul.py}，或任意使用
\texttt{T.gemm} 的 kernel，分别以 \texttt{sm\_90a} 和 \texttt{sm\_100a}
为 target 编译。 本题只要求编译，不需要实际使用 H 卡。

保存生成的 CUDA 源码与 lowering 输出，并填写：

\begin{longtable}[]{@{}lll@{}}
\toprule\noalign{}
& sm\_90a & sm\_100a \\
\midrule\noalign{}
\endhead
\bottomrule\noalign{}
\endlastfoot
选中的 Tensor Core 指令 & & \\
descriptor 在哪里、由谁生成 & & \\
smem swizzle 布局在哪一步确定 & & \\
数据由谁搬入 smem & & \\
\end{longtable}

与 M2--M4 中的手写实现进行对照，并回答：

\begin{enumerate}
\def\labelenumi{(\alph{enumi})}
\item
  哪些硬件相关的决策已经由 DSL 自动完成？
\item
  哪些参数仍然需要程序员决定，例如 tile 尺寸和 stages？
\end{enumerate}

最后在 assignment01 7.5 的``谁负责''表中补充一行：

\texttt{Tensor\ Core\ 指令选择与供数布局}

\section*{团队选做（推荐）}\label{ux56e2ux961fux9009ux505aux63a8ux8350}
\addcontentsline{toc}{section}{团队选做（推荐）}

2--4 人一组，从 \texttt{team/} 中的两个题目任选一个，也可以全部完成：

\begin{itemize}
\tightlist
\item
  C1：FlashKDA 官方 kernel 当前使用 SM80 MMA，分析迁移到 SM100
  是否值得；
\item
  C2：MiniMax M3 MSA 的小 batch decode 当前使用
  Triton，分析是否值得实现专用 kernel。
\end{itemize}

每道题分为三个阶段：

\texttt{复现与测量\ →\ 分析（结论\ +\ 证据）→\ 挑战}

挑战阶段不要求一定得到正向加速。如果实验结果表明继续优化收益有限，
只要能够用数据说明为什么不值得继续做，同样视为有效结论。

答辩时间为 10 分钟，另有 5 分钟提问；提问优先由选择另一道团队题的
小组提出。具体要求见 \texttt{team/README.md}。

团队题不限制 AI 工具的使用。

\section*{提交}\label{ux63d0ux4ea4}
\addcontentsline{toc}{section}{提交}

\begin{itemize}
\tightlist
\item
  \textbf{代码}：提交所有动手题的实现与判测输出；FROM-SCRATCH
  题同时保留判测脚本的 PASS 记录。
\item
  \textbf{报告}：包含纸面题解答、实验表格与性能归因，以及 DEBUG
  题的现象记录和修改说明。所有实验数据注明使用的 GPU。
\item
  \textbf{团队题}：提交代码、报告并完成答辩，具体要求见
  \texttt{team/README.md}。
\end{itemize}

\end{document}
